\begin{abstract}
% v10. New abstract supplied by the author (conventional context->prior work->gap->method->
% findings->implication arc). One correctness fix vs the pasted text: "cosmetic edits" -> "restate
% or decline ... leaving outputs effectively unchanged", to match the paper's classifier, which
% defines the non-genuine class as meta-responses (declines/restatements), not small edits.
% No em-dashes. Numbers: 39/13, 0.74, 1.16.
Large language models are increasingly asked to revise their own outputs over multiple conversation
turns, often by users who cannot specify what should change. Prior work on self-correction has
largely examined single rounds of self-critique on tasks with verifiable answers, finding that models
improve mainly when external feedback is available. Whether repeated, undirected revision requests
help or harm output quality remains unexamined. We study this setting directly, running five-turn
revision conversations with six LLMs on 40 tasks across five domains, from code and data reasoning to
business and creative writing. Under undirected requests, genuine revision declines from 39\% of
responses at the second turn to 13\% by the fifth, with models instead restating or declining in ways
that leave outputs effectively unchanged. Revisions that do occur reduce quality by 0.74 levels on a
five-point scale, and every model's first draft rates highest. A single targeted critique reverses
this pattern, improving quality by 1.16 levels. These results suggest that intrinsic self-correction
fails not from lack of capacity but from lack of direction: models revise effectively only when told
what is wrong.
\end{abstract}
